\documentclass[11pt,letterpaper]{article}
\usepackage[margin=1in]{geometry}
\usepackage{xcolor}
\usepackage{hyperref}
\usepackage{enumitem}
\usepackage{titlesec}
\usepackage{fancyhdr}
\usepackage{parskip}
\usepackage{microtype}
\usepackage{xurl}

\hypersetup{colorlinks=true, linkcolor=blue, urlcolor=blue, citecolor=blue}

% Aggressive line-breaking inside \url{} / \nolinkurl{} for HuggingFace-style identifiers.
\renewcommand{\UrlBreaks}{\do\/\do\-\do\_\do\.\do\&\do\?\do\=}

% Allow LaTeX to stretch line spacing slightly rather than overrunning the margin.
\sloppy
\emergencystretch=3em
\hyphenpenalty=100
\tolerance=2000

% Reviewer-comment block (italic, blue) vs author-response block (regular, black).
\newcommand{\reviewer}[1]{\par\medskip\textcolor{blue!75!black}{\textit{#1}}\par\medskip}
\newcommand{\response}[1]{\par#1\par\medskip}
% Quoted excerpt of the revised manuscript, so reviewers need not consult the
% paper. Arg 1 = location label; arg 2 = quoted text (newly added words in blue).
\definecolor{msgray}{gray}{0.28}
\newcommand{\fromms}[2]{\par\smallskip{\small\setlength{\leftskip}{1.4em}\setlength{\rightskip}{0.6em}%
\textcolor{msgray}{\textbf{$\hookrightarrow$ Revised manuscript, #1:}}\par
\textcolor{msgray}{\itshape #2}\par}\smallskip}

\titleformat{\section}{\Large\bfseries\color{black}}{\thesection}{0.5em}{}
\titleformat{\subsection}{\large\bfseries\color{black!85}}{\thesubsection}{0.5em}{}

\pagestyle{fancy}
\fancyhf{}
\fancyhead[L]{\textbf{CLINES — Response to Reviewers}}
\fancyhead[R]{Manuscript \texttt{bmjdh-2026-000027}}
\fancyfoot[C]{\thepage}

\title{\vspace{-2em}\textbf{Response to Reviewers}\\[0.4em]\large CLINES: Clinical LLM-based Information Extraction and Structuring Agent\\[0.2em]\normalsize BMJ Digital Health \& AI --- Manuscript \texttt{bmjdh-2026-000027}}
\author{}
\date{May 29, 2026}

\begin{document}
\maketitle
\thispagestyle{fancy}
\vspace{-3em}

\section*{Overview}

We thank the Editor and the five reviewers for their detailed and constructive feedback. The revised manuscript incorporates substantial additions in response to the major concerns raised, including:

\begin{itemize}[leftmargin=*]
    \item A structured \textbf{inter-annotator agreement (IAA)} study with cross-annotation on three notes per dataset (4CE, CORAL), reporting mention-level F1 (0.82), assertion Cohen's~$\kappa$ (0.85), and per-attribute exact-match rates (Methods (Inter-annotator agreement); Figure~S1A; Supplementary Table~S5).
    \item \textbf{Bootstrap 95\% confidence intervals} and paired-bootstrap, FDR-corrected pairwise significance tests on all primary F1 comparisons (the Results; Figure~4A).
    \item A \textbf{seven-class hallucination taxonomy} with stratified manual judging of 700 candidate false positives, extrapolated population-level rates, and full case studies (Methods; the Results; Figure~S1C-D).
    \item \textbf{Cost, latency, and computational efficiency} reporting per dataset, per model, and per LLM-invocation type (the Results; Figure~4B).
    \item \textbf{Expanded baselines}: o3-mini single-prompt and GPT-4o chain-of-thought conditions (addressing the previous baseline asymmetry), plus two off-the-shelf clinical deep-learning encoders (BERT-base clinical NER and GatorTron-base) for direct LLM-vs-DL comparison (Results; Figure~4D).
    \item \textbf{Component-wise ablation} of the four pipeline stages (SapBERT normalization, semantic chunking, date module, cross-chunk reconciliation; Figure~4C).
    \item \textbf{MI-CLAIM and TRIPOD+AI} reporting checklists in place of STROBE (supplementary materials).
    \item Multiple smaller clarifications regarding annotation protocol, data governance, computational requirements, and downstream interoperability.
\end{itemize}

All changes in the revised manuscript are marked in \textcolor{blue}{blue}. Cross-references in this document point to the revised manuscript page/section numbers.

\hrulefill

\section{Reviewer 1}

\subsection*{Major concerns}

\subsubsection*{R1.1 --- Missing Inter-Annotator Agreement}
\reviewer{``The absence of inter-annotator agreement (IAA) reporting is the most significant methodological weakness\dots\ I would strongly encourage the authors to either report IAA for existing annotations or double-annotate a representative subset and report agreement statistics (e.g., Cohen's kappa, F1-based IAA) in a revised manuscript.''}
\response{We thank the reviewer for this critical observation. We have conducted a structured cross-annotation study on three notes per dataset (4CE $n=3$; CORAL $n=3$; total 2{,}675 candidate-mention rows). For each note, a second annotator, blinded to the original assignment, independently re-annotated the AI-suggested draft. We report mention-level agreement as set-based F1 (the standard for clinical-NER IAA; Uzuner et al.\ \textit{JAMIA} 2010 17:519--523), assertion agreement as Cohen's~$\kappa$ on the 5-class collapsed label space (kept-by-both subset), value and unit as exact-match rates over all kept-by-both rows, and begin- and end-date as exact-match rates restricted to rows where both annotators recorded a date. After applying pre-specified structural-disambiguation guidelines, pooled mention F1 was \textbf{0.82}, comparable to the pair-wise F1 inter-annotator agreement (0.81--0.88) reported for the i2b2 medication-challenge community annotation experiment (Uzuner et al.\ \textit{JAMIA} 2010 17:519--523); pooled assertion $\kappa = 0.85$, in the substantial-agreement range (Landis \& Koch 1977). Pooled value and unit exact-match rates were 0.71 and 0.98, and pooled begin- and end-date exact-match rates were 0.92 and 0.94. For MIMIC-III, where each paragraph was annotated by a single expert, we explicitly reframe F1 as agreement-with-annotator. The study design appears in the revised Methods (``Inter-annotator agreement''); the headline numbers appear in the Results (``Expanded baselines, cost, and ablation'') and in the Discussion (Limitation iv); the per-note breakdown, the five structural-disambiguation patterns, and the comparison to the Uzuner 2010 benchmark are in Supplementary Materials (``Inter-Annotator Agreement: Methodology and Per-Note Detail'').}
\fromms{Supplementary Materials, ``Inter-Annotator Agreement: Methodology and Per-Note Detail'' (Headline agreement)}{``\textcolor{blue}{Across the cross-annotated subset, pooled mention F1 was \textbf{0.82}, comparable to the pair-wise F1 inter-annotator agreement (0.81--0.88) reported for the i2b2 medication-challenge community annotation experiment; assertion $\kappa = 0.85$ (substantial agreement); value and unit exact-match rates 0.71 and 0.98; begin- and end-date exact-match rates 0.92 and 0.94.}''}

\subsubsection*{R1.2 --- Confidence Intervals and Statistical Testing on Core Results}
\reviewer{``The main results (Figure 3A-D) report point-estimate F1 scores without confidence intervals or formal statistical comparisons\dots\ Bootstrap confidence intervals and/or permutation tests should be provided for all primary comparisons.''}
\response{We now report \textbf{percentile bootstrap 95\% confidence intervals} ($B=2000$ document-level resamples) for every primary F1 comparison. To keep the main-text Figure~3A--D legible we report the per-corpus point F1 there as in the original submission, and tabulate the full per-field point estimates with 95\% bootstrap CIs in Supplementary Materials (``Detailed Performance Metrics with Bootstrap Confidence Intervals''); Figure~3E additionally renders the length-quantile CLINES-vs-baseline comparison with 95\% shaded bootstrap ribbons. For example, the reference CLINES (o3-mini-medium) code-level F1 was 0.874 [95\% CI 0.849--0.896] on 4CE, 0.848 [0.834--0.865] on CORAL--Pancreas, and 0.814 [0.787--0.838] on CORAL--Breast. Pairwise significance was assessed with a \textbf{paired bootstrap} on the shared document resamples---which yields a confidence interval on the per-document F1 \emph{difference} directly and is at least as powerful as a label-permutation test for this paired design---with Benjamini--Hochberg FDR correction within each metric family (per-field bootstrap CIs are tabulated in the Supplementary Material). CLINES (o3-mini-medium) exceeded both backbone alternatives on code, assertion, value, and unit extraction ($\Delta$F1 $+0.01$ to $+0.21$ versus GPT-4o and Llama-3.1-405B). The Code-field gains and most assertion/value cells were significant at $p_{\mathrm{FDR}} < 0.05$; the smallest gaps (e.g.\ 4CE Unit at $+0.01$) fell within bootstrap noise and were not significant — we note this explicitly in the revised Results. In the interest of honest reporting we note that the date column was \emph{not} reliably distinguishable among the strongest backbones at this sample size, and we say so explicitly in the revised Results.}

\subsubsection*{R1.3 --- Small and Heterogeneous Evaluation Sets}
\reviewer{``The annotated evaluation sets are notably small\dots\ stability analyses (e.g., bootstrap variance of F1 across document-level resamples) to characterise how sensitive the reported metrics are to the specific documents in each evaluation set.''}
\response{We added a document-level resampling stability analysis ($B=2000$ document-level bootstrap resamples for every model$\times$dataset$\times$task cell; Figure~3E). This quantifies precisely the concern raised: on the smallest set (CORAL--Breast, $n=13$) the bootstrap standard deviation of F1 reaches $\approx 0.04$--0.05 and CI half-widths approach $\pm 0.05$--0.10, so F1 differences smaller than $\approx 0.10$ are not statistically resolvable at this sample size. We therefore reframed the Results to rest the principal claims on comparisons that survive this uncertainty (CLINES's advantage on code/assertion/value/unit, whose effect sizes lie well above the resolvable threshold) rather than on the precise magnitude of any single F1 estimate, and we state the limited resolution explicitly in the revised Discussion. Figure~3E additionally annotates per-bin note counts so that small, unstable bins are not over-interpreted (cf.\ R5.3.5).}

\subsubsection*{R1.4 --- Inappropriate Reporting Guideline (STROBE)}
\reviewer{``The STROBE\dots\ framework is not fit for purpose here\dots\ The authors should instead adopt a reporting guideline appropriate for AI and machine learning studies in healthcare, such as MI-CLAIM, TRIPOD+AI, or STARD-AI.''}
\response{We agree. The STROBE checklist has been removed and replaced with a \textbf{MI-CLAIM} checklist (Norgeot et al., Nat Med 2020) and a \textbf{TRIPOD+AI} checklist (Collins et al., BMJ 2024), both included in the supplementary materials. STARD-AI was considered but is intended for diagnostic-accuracy studies and is not applicable to information-extraction benchmarks.}

\subsubsection*{R1.5 --- No Cost, Latency, or Computational Efficiency Reporting}
\reviewer{``\dots\ the absence of any cost or latency analysis is a significant omission\dots\ I would encourage the authors to provide a transparent accounting of the total number of LLM invocations per note, stratified by model and dataset, along with per-note wall-clock processing time and estimated API cost.''}
\response{We have added a full cost/latency/invocation breakdown. The main-text Methods (``Implementation and LLM configuration'') retains the original per-chunk token statistics and the main-text Results (``Expanded baselines, cost, and ablation'') summarizes per-note cost across deployments; the full per-dataset, per-backbone breakdown---LLM invocations per note, wall-clock seconds, prompt + completion tokens, and USD cost at Azure OpenAI list rates---is in Supplementary Materials (``Extended Results: Cost, latency, and computational efficiency'') and Figure~4B. We acknowledge that single-note inference involves more LLM calls than the ``four-step'' architectural framing implies (cf.\ R1.12 below); the cost accounting reflects the actual call count.}
\fromms{Supplementary Materials, ``Methodology Extensions'' (Llama-3.1-405B serving configuration and architectural-stage taxonomy)}{``\textcolor{blue}{We distinguish the four \emph{architectural stages} (chunk $\rightarrow$ extract $\rightarrow$ normalize $\rightarrow$ reconcile) from the per-chunk LLM invocations: each chunk triggers multiple invocations across the extract/normalize/date sub-prompts, so the per-note call count exceeds four. Full per-step invocation, token, time, and cost breakdown appears with the cost analysis in Figure~4B.}''}

\subsection*{Moderate concerns}

\subsubsection*{R1.6 --- ``Agentic'' Overclaim}
\reviewer{``In the current AI literature, agentic systems are typically characterised by autonomous decision-making\dots\ characterising CLINES as `agentic' overstates the architectural sophistication.''}
\response{We agree the term ``agentic'' carries connotations that exceed the actual architectural sophistication of CLINES. Throughout the revised manuscript, ``agentic pipeline'' has been replaced by ``\textbf{multi-step modular pipeline}'' (or ``modular orchestration'' where flow is the salient property). The title retains the original acronym for traceability but we no longer claim agentic autonomy in the architectural descriptions.}
\fromms{Discussion, ``Positioning versus existing approaches''}{``In CLINES, a \textcolor{blue}{multi-step} orchestration---chunking, tool-guided extraction, normalization, and date reconciliation---delivers structured outputs with transparent intermediate artifacts and maintains robustness across diverse note types.'' (The term ``agentic'' is likewise replaced by ``multi-step'' in the Dataset-level variation and Conclusion subsections.)}

\subsubsection*{R1.7 --- Incomplete Baseline Landscape}
\reviewer{``\dots\ no comparison\dots\ against GPT-4o or o3-mini using chain-of-thought prompting in a single call\dots, which would be a more informative ablation of the multi-step decomposition strategy.''}
\response{We have added two additional baseline conditions (the Results; Figure~4D): (i) \textbf{o3-mini single-prompt} matching the previously-reported GPT-4o single-prompt setting (closing the asymmetry that R5 flagged separately at point 4.1), and (ii) \textbf{GPT-4o chain-of-thought single-prompt} with explicit \emph{Let's think step by step} instruction plus the structured output schema. These additions disentangle backbone-model contribution from pipeline-architecture contribution. On code-level extraction the single-call baselines fell far below CLINES: o3-mini single-prompt reached F1 0.50 / 0.43 / 0.36 on 4CE / CORAL--Pancreas / CORAL--Breast and GPT-4o chain-of-thought 0.40 / 0.34 / 0.27, versus CLINES 0.87 / 0.85 / 0.81---an architecture-attributable gain of $+0.37$ to $+0.54$ F1 that cannot be explained by backbone model or prompt style alone (Figure~4D). We additionally include off-the-shelf full-size clinical-transformer baselines (cf.\ R3.8); see R3.8 response.}

\subsubsection*{R1.8 --- Prompt Sensitivity and Ablation}
\reviewer{``No ablation study is provided to quantify the contribution of individual pipeline stages\dots\ The authors acknowledge that prompts were not optimised per dataset or model but provide no sensitivity analysis.''}
\response{We added a component-wise ablation study (Figure~4C), run on two backbones (GPT-4o-1120 and GPT-4o-mini). Each of the four major pipeline components was independently disabled and code-level F1 re-measured relative to the full pipeline: (a) SapBERT-based normalization replaced by LLM-only normalization; (b) semantic chunking replaced by fixed-length chunking; (c) date module disabled; (d) cross-chunk reconciliation disabled (Step~4). Removing SapBERT normalization caused by far the largest degradation (mean $\Delta$code-F1 $-0.447$ with GPT-4o and $-0.414$ with GPT-4o-mini), confirming UMLS normalization as the load-bearing component. Disabling cross-chunk reconciliation cost $-0.078$ (GPT-4o) / $-0.221$ (GPT-4o-mini), while semantic chunking ($-0.021$ / $-0.031$) and the date module ($-0.003$ / $-0.008$) contributed smaller amounts. The reconciliation step mattered \emph{more} for the cheaper backbone, indicating the deterministic scaffolding partly compensates for a weaker model. The ablation used $n=2$ representative notes per dataset and reports point estimates (noted as a limitation).}

\subsubsection*{R1.9 --- Hallucination Assessment}
\reviewer{``No systematic quantification of hallucination rates is provided\dots\ A structured error analysis categorising false positives by error type (hallucinated entity, incorrect normalization, wrong assertion, fabricated date) would significantly strengthen the evaluation.''}
\response{We implemented a stratified manual-judging study of 700 candidate false-positive predictions under a seven-class hallucination taxonomy: \texttt{fabricated\_entity}, \texttt{span\_boundary\_error}, \texttt{wrong\_assertion}, \texttt{wrong\_code}, \texttt{wrong\_value}, \texttt{wrong\_date}, and \texttt{not\_an\_error} (annotation gap / synonym equivalence / guideline ambiguity). Extrapolating to the full FP population ($N=7{,}171$) via post-stratified Horvitz-Thompson estimation, we find that 54.3\% of apparent FPs are \texttt{not\_an\_error} (annotation gap), and each of the six true-hallucination categories accounts for $\leq 17\%$ of FPs --- with \texttt{wrong\_code} (16.7\%) being the most common true-hallucination category. See revised Methods, the Results, Figure~S1C-D, and full case studies in the Supplementary Material (Hallucination Taxonomy and Case Studies).}

\subsubsection*{R1.10 --- Reference Error (Uzuner 2011)}
\reviewer{``Reference 14 is cited in the context of the `i2b2-2010 concepts' benchmark, but the reference itself is Kohane et al.\ 2012\dots\ The correct citation should be Uzuner et al.\ (2011).''}
\response{Corrected. The Uzuner et al.\ 2011 reference (\textit{JAMIA} 18(5):552--556) replaces the previous citation in both the Introduction (where transformer encoders are reported to perform well on ``i2b2 concept extraction'') and the Methods Baselines paragraph (where Clinical-MobileBERT/DistilBERT are described as ``fine-tuned on i2b2-2010 concepts''); the original Kohane et al.\ 2012 entry has been removed from the reference list because no remaining main-text claim requires it.}

\subsubsection*{R1.11 --- Prompt Examples Redacted from Supplementary Material}
\reviewer{``\dots\ all in-context few-shot examples have been replaced with placeholders\dots\ Without the actual examples, an independent researcher cannot reconstruct the pipeline's behaviour from the supplement alone.''}
\response{Full few-shot examples have been restored in the revised supplementary material. We also confirm that the complete prompt suite, including all few-shot examples, will be published in the public GitHub repository associated with this work (URL specified in the revised Data Availability statement).}

\subsubsection*{R1.12 --- Discrepancy Between Architecture and Pipeline Complexity}
\reviewer{``The main manuscript describes CLINES as a `four-step pipeline.' However, the supplementary prompt suite reveals at least 11 distinct prompt types, implying that a single chunk traversal involves considerably more LLM invocations than the four-step framing suggests.''}
\response{Acknowledged and clarified. The revised Methods now explicitly distinguishes between the four \emph{architectural stages} (chunk $\rightarrow$ extract $\rightarrow$ normalize $\rightarrow$ reconcile) and the underlying \emph{LLM invocations}, which include attribute extraction, assertion classification, date normalization, JSON repair, and other helper sub-calls. The total invocation count per note is now reported per dataset and per model in Figure~4B (cf.\ R1.5 response).}

\subsubsection*{R1.13 --- Unevaluated Relation Extraction Module}
\reviewer{``Supplementary Section 2.2 describes a complete entity relationship linking module\dots\ Yet the main manuscript does not evaluate relation extraction performance at all.''}
\response{The relation extraction module was deployed but not evaluated in the original submission because gold-standard relation annotations are not available for any of the four corpora. The revised Methods explicitly flags this scope limitation; the relation extraction prompt is retained as an architectural component of the deployed system but its quantitative evaluation is deferred to future work where annotated relation data are available.}

\subsubsection*{R1.14 --- UMLS Retrieval Parameters Underspecified}
\reviewer{``The UMLS retrieval procedure (Supplementary Algorithm 1) is described at a very high level\dots\ which SapBERT checkpoint\dots\ FAISS index type\dots\ top-$k$\dots\ similarity threshold\dots\ UMLS version\dots''}
\response{The revised supplementary material now specifies the retrieval procedure exactly as implemented: SapBERT checkpoint \nolinkurl{cambridgeltl/SapBERT-from-PubMedBERT-fulltext}; L2-normalized embeddings searched with FAISS using an exact inner-product index (\texttt{IndexFlatIP}) for concept lists below 50{,}000 terms and an inverted-file index (\texttt{IndexIVFFlat}, \texttt{nlist}~$=\min(4096,\max(N/50,256))$, \texttt{nprobe}~$=32$, inner-product metric) for larger lists; candidates are ranked by cosine similarity and the top-ranked candidates are passed to an LLM-assisted disambiguation step that selects the final concept, rather than applying a fixed similarity cutoff. The UMLS subset, release version, and source-vocabulary filter used to build the concept index are stated in revised Supplementary Material (UMLS retrieval).}

\subsubsection*{R1.15 --- Research Ethics: Annotator Credentialing and Commercial API Data Governance}
\reviewer{``The annotation process involved five human annotators reviewing de-identified clinical notes from credentialed-access datasets. The manuscript does not confirm whether all annotators held appropriate data access credentials\dots\ Additionally, the pipeline transmits de-identified clinical text to commercial API endpoints (OpenAI)\dots''}
\response{Two clarifications added to the revised manuscript:
\begin{itemize}[leftmargin=*]
    \item \textbf{Annotator credentialing}: Methods (Annotation protocol) and the Ethics Approval section now explicitly state that all annotators completed CITI Human Subjects Research training and held appropriate credentials for the datasets they accessed (PhysioNet credentialed access for MIMIC-III and CORAL; 4CE consortium credentialing for 4CE data).
    \item \textbf{API governance}: All hosted-API LLM inference (GPT-4o, o3-mini) was performed via HMS Information Technology's managed Azure AI service (HMS Azure AI; \url{https://it.hms.harvard.edu/service/azure-ai}), an HMS-administered Microsoft Azure OpenAI deployment scoped to Harvard data-security Level~4 projects. Only de-identified clinical text was processed (consistent with HMS Azure AI's exclusion of HIPAA-regulated data), and Microsoft Azure OpenAI Service does not retain submitted prompts/completions for model training. This deployment configuration was reviewed for compliance with the MIMIC-III, 4CE, and CORAL DUAs prior to any inference. See Methods (Implementation and LLM configuration).
\end{itemize}}
\fromms{Methods, ``Implementation and LLM configuration'' (brief) and Supplementary Materials, ``Methodology Extensions'' (Compute and governance configuration; full)}{``\textcolor{blue}{Hosted-API inference (GPT-4o, o3-mini) was performed via HMS Information Technology's managed Azure AI service (HMS Azure AI) on de-identified clinical text, with no data retained or used to train provider models. The open-weight Llama-3.1-405B backbone was served on-premises. All annotators completed CITI Human Subjects Research training and held the access credentials required by their respective institutions for the de-identified study data.}''}

\subsubsection*{R1.16 --- Extraction of Sensitive Demographic Attributes}
\reviewer{``\dots\ the pipeline includes a dedicated prompt (Supplement 2.6) that routinely extracts race, ethnicity, gender, and zip code\dots\ The manuscript does not discuss the intended use cases\dots\ whether extraction performance varies across demographic subgroups, or the potential risks\dots''}
\response{We added an ``Ethical considerations in demographic-attribute extraction'' subsection to Supplementary Materials (``Deployment, Interoperability, and Ethical Considerations''), with a brief pointer in the main-text Discussion (``Deployment, interoperability, and ethical considerations''). The added text: (i) clarifies that demographic extraction is included to support downstream cohort discovery and equity audits, not for any model-personalization use case; (ii) acknowledges that we have not evaluated subgroup-specific extraction performance, and explicitly recommends this as a required pre-deployment evaluation step for any institutional adopter; and (iii) discusses the responsibility of downstream users to apply appropriate access controls when demographic fields are stored.}

\subsection*{Minor issues}

\subsubsection*{R1.17 --- MIMIC-III Entity F1 lower (0.69) than other datasets}
\reviewer{``\dots\ MIMIC-III (0.69 for the best CLINES variant) is notably lower than on other datasets (0.81-0.87)\dots\ heterogeneity of critical care notes, annotation inconsistencies\dots, or limitations of the extraction prompts for ICU-specific terminology.''}
\response{Discussed in the revised Discussion (``Dataset-level variation and generalization''). We attribute the MIMIC-III gap primarily to (a) the single-annotator-per-paragraph protocol used for MIMIC-III, which raises the variability of the gold standard relative to the consensus-annotated 4CE and CORAL data (cf.\ R1.1, R5.4.2), and (b) ICU-specific terminology that is underrepresented in the few-shot examples. We explicitly recommend dataset-specific few-shot expansion as a deployment-stage tuning step.}

\subsubsection*{R1.18 --- Abstract Structure (BMJ style)}
\reviewer{``The abstract uses Lancet-style headings (Background, Methods, Findings, Interpretation) rather than BMJ-style structured abstract headings.''}
\response{Restructured. The revised abstract uses BMJ Digital Health \& AI's required headings: \textbf{Objective}, \textbf{Methods and Analysis}, \textbf{Results}, \textbf{Conclusion} (cf.\ R5.2.1). A summary box (``What is already known / What this study adds / How this study may affect research, practice, or policy'') has also been added immediately after the abstract (cf.\ R5.2.2).}

\subsubsection*{R1.19 --- STROBE Checklist Text Artifacts}
\reviewer{``Several cells in the STROBE checklist (pages 39-42) contain garbled text\dots''}
\response{The STROBE checklist has been removed entirely (cf.\ R1.4) and replaced with cleanly-typeset MI-CLAIM and TRIPOD+AI checklists.}

\subsubsection*{R1.20 --- Missing PPI Statement}
\reviewer{``BMJ journals typically require a Patient and Public Involvement (PPI) statement\dots''}
\response{Added. The revised manuscript includes a Patient and Public Involvement statement immediately before the Ethics Approval section, confirming that no patients or members of the public were involved in the design, conduct, reporting, or dissemination of this study, given that it analyzes only fully de-identified secondary data.}

\subsubsection*{R1.21 --- IRB Authority}
\reviewer{``The ethics approval section references `the policies of our institution' without naming the specific institution or IRB\dots''}
\response{Clarified. The vague phrase ``the policies of our institution'' has been replaced with a specific and accurate statement. This study analyzes only fully de-identified, publicly available or credentialed-access secondary datasets (MIMIC-III, CORAL, 4CE), each used under its own data-use agreement. Under the U.S. Common Rule (45~CFR~46.104), secondary analysis of de-identified data that are not readily identifiable does not constitute human-subjects research and therefore did not require Institutional Review Board approval or a formal exemption application at the authors' institution (Harvard Medical School). All annotators completed CITI Human Subjects Research training and held the access credentials required for each dataset. See the revised Ethics Approval section.}
\fromms{Ethics Approval}{``\textcolor{blue}{Because the analysis involves only fully de-identified secondary data that are not readily identifiable, it does not constitute human-subjects research under the U.S. Common Rule (45~CFR~46.104) and did not require Institutional Review Board approval or a formal exemption determination at the authors' institution (Harvard Medical School)\dots\ All annotators completed CITI Human Subjects Research training and held the access credentials required for each dataset prior to data access.}''}

\subsubsection*{R1.22 --- Terminology Consistency (value/unit)}
\reviewer{``The manuscript inconsistently uses `value\&unit' (abstract), `value-unit' (results), and `value/unit' (discussion)\dots''}
\response{Unified to \texttt{value+unit} throughout, in line with the schema used in the deployed pipeline.}

\hrulefill

\section{Reviewer 2}

\subsection*{Major comments}

\subsubsection*{R2.1 --- ``Zero-shot'' naming}
\reviewer{``\dots\ the authors clearly incorporates prompt engineering, task decomposition, retrieval-augmented reasoning (e.g., SapBERT), and rule-based post-processing. Under standard NLP definitions, this does not constitute true zero-shot learning. Please reconsider naming the `zero-shot' evaluation.''}
\response{Agreed. ``Zero-shot evaluation'' has been replaced throughout by ``\textbf{no task-specific supervised training}'' or ``\textbf{evaluation without fine-tuning on the target annotation distribution}''. The revised Methods explicitly states that the pipeline relies on prompt engineering with few-shot in-context examples, retrieval-augmented normalization, and post-processing rules, none of which qualify as zero-shot in the strict NLP sense.}
\fromms{Methods, ``Evaluation metrics and protocol''}{``In the \textcolor{blue}{training-free setting (prompt engineering with few-shot in-context examples, retrieval-augmented normalization, and post-processing rules; not zero-shot in the strict sense)}, we evaluated\dots''}

\subsubsection*{R2.2 --- Annotation reliability}
\response{See R1.1 response. IAA cross-annotation study added.}

\subsubsection*{R2.3 --- Hallucination measurement}
\response{See R1.9 response. Seven-class hallucination taxonomy added.}

\subsubsection*{R2.4 --- Generalizability claims}
\reviewer{``\dots\ annotated data are unevenly distributed across datasets, not covering several common note types (e.g., discharge summaries, operative or radiology reports). Claims regarding generalizability should be framed accordingly.''}
\response{We have tempered the generalizability framing throughout the revised manuscript. The Discussion now explicitly enumerates note types \emph{not} covered (discharge summaries, operative notes, radiology reports) and recommends each adopter to validate CLINES on local note-type distributions before clinical deployment.}

\subsection*{Minor comment}

\subsubsection*{R2.5 --- Tone down ``replace chart-review''}
\reviewer{``Would suggest toning down absolute phrasing, indicating that CLINES can replace chart-review with hallucination still present, and would still need manual reviews''}
\response{Adjusted. The Discussion now states ``CLINES can \emph{substantially reduce} manual chart review for entities, assertions, value+unit, and dates, while retaining a verifiable audit trail. We recommend that clinical adopters maintain human review for any extraction destined for use in safety-critical decisions, given the residual hallucination rates quantified in the Results.''}

\hrulefill

\section{Reviewer 3}

\subsubsection*{R3.1 --- Additional citations in Introduction (EF, tumor size)}
\reviewer{``\dots\ measurements such as cardiac ejection fraction or tumor size often do not fit neatly into predefined problem lists or billing codes\dots\ should be supported with appropriate references.''}
\response{Citations added. The opening Introduction sentence on measurements (e.g., ejection fraction, tumor size) that escape structured problem lists and billing codes is now supported by two references on the limitations of structured EHR capture for such free-text clinical detail (Tayefi et al., \emph{WIREs Comput Stat} 2021; Sushil et al., 2024).}
\fromms{Introduction (para.\ 1)}{``\dots measurements like cardiac ejection fraction values or tumor size often do not fit neatly into predefined problem lists and billing codes yet are central to clinical interpretation\,\textcolor{blue}{[Tayefi 2021; Sushil 2024]}.''}

\subsubsection*{R3.2 --- Intermediate ML/DL work in Introduction}
\reviewer{``The literature framing in the Introduction moves directly from rule-based approaches to transformer-based LLM methods\dots\ what ML/DL approaches have previously been applied to this problem, and what limitations of those methods motivated the current approach?''}
\response{A new sentence has been inserted into the Introduction surveying the intermediate neural era between rule-based systems and modern LLMs: BiLSTM-CRF sequence taggers (Lample et al.\ 2016) and, subsequently, domain-pretrained transformer encoders BioBERT (Lee et al.\ 2020) and ClinicalBERT (Alsentzer et al.\ 2019). The limitations motivating LLM-based methods --- requirement for sizeable task-specific labeled corpora, brittleness to phrasing variation and institutional shift, and treatment of entity detection in isolation rather than joint normalization, attribute, and temporal extraction --- are explicitly enumerated. (GatorTron and Clinical-MobileBERT remain cited in the subsequent paragraph as domain-scaled encoders.)}
\fromms{Introduction (para.\ 3)}{``\textcolor{blue}{Between these rule-based systems and modern large language models, a generation of supervised neural methods advanced clinical entity recognition and normalization---BiLSTM-CRF sequence taggers and, subsequently, domain-pretrained transformer encoders such as BioBERT and ClinicalBERT [Lample 2016; Lee 2020; Alsentzer 2019]. These improved span detection but require sizeable task-specific labeled corpora, remain brittle to phrasing variation and institutional shift, and typically address entity detection in isolation rather than joint normalization, attribute, and temporal extraction.}''}

\subsubsection*{R3.3 --- Step 1 chunking detail}
\reviewer{``\dots\ Step 1, the chunking strategy is not sufficiently explained. It remains unclear whether a sliding-window approach was used\dots''}
\response{Methods (Note chunking) has been expanded. The chunking algorithm is now described as: TikToken (cl100k\_base encoding) for token counting, SemChunk for boundary selection at sentence and paragraph breaks, target chunk size $\leq 768$ tokens, \textbf{no sliding window} (chunks are non-overlapping), \textbf{no truncation} of source text. Core metadata (admission date, discharge date, patient identifier) is replicated in every chunk header to preserve temporal alignment. The justification for $768$ tokens is included (cf.\ R5.3.4 on SemChunk capacity).}
\fromms{Supplementary Materials, ``Methodology Extensions'' (Semantic chunking and overlap policy); pointer in Methods, ``Step 1 -- Note chunking''}{``\textcolor{blue}{Chunks are non-overlapping (no sliding window is used) and no source text is truncated; the segmenter operates on the full note. Token budgets and the recommended sentence-aware splitter are inherited from the \texttt{semchunk} library; the tokeniser used to count budget is the model's native tokeniser. Within each chunk the prompt operates on the full sentence boundary set so that adjacent mentions necessary for assertion or value inference remain co-resident.}''}

\subsubsection*{R3.4 --- Step 4 reconciliation detail}
\response{See R3.3 plus expanded text on Step~4 in Methods. Cross-chunk reconciliation is described as: (i)~chunk-local entity tags are re-indexed into a document-global namespace via a running tag offset so that downstream per-step attribute/date outputs can be joined back to a single global entity table; (ii)~per-step attribute and date outputs are joined to entity records by global tag, preserving each chunk's locally-extracted attribute values; and (iii)~relative-date expressions are resolved against the note's admission timestamp. Because chunks are non-overlapping (R3.3) a given character span appears in at most one chunk, so explicit cross-chunk span deduplication is unnecessary; same-surface mentions arising in different chunks are retained as separate records, each carrying their chunk-local context's attribute assignments (a deliberate choice to preserve context-conditional differences such as a finding being \emph{Present} in one section and \emph{Historical} in another).}

\subsubsection*{R3.5 --- Ablation study}
\response{See R1.8 response. Component-wise ablation added (Figure~4C).}

\subsubsection*{R3.6 --- ``Words'' definition}
\reviewer{``\dots\ unclear whether `words' refers to whitespace-delimited items, tokenizer-based units, or another operational definition.''}
\response{Defined explicitly in the revised Methods: ``\emph{words} = whitespace-and-punctuation-delimited tokens; stopwords are included in word counts; numerical values and abbreviations are counted as single words.''}

\subsubsection*{R3.7 --- ``Clinical phrases'' definition}
\reviewer{``A similar issue applies to the use of the term `clinical phrases.'\dots''}
\response{Defined in the revised Methods: ``\emph{clinical phrases} = contiguous spans corresponding to a single UMLS-mappable concept (or a single numerical-value-plus-unit), as marked by the annotators.''}

\subsubsection*{R3.8 --- Missing prior DL baselines}
\reviewer{``The baseline comparison in the Methods appears limited in that it does not include prior deep learning-based approaches.''}
\response{We added two full-size, publicly available clinical deep-learning baselines evaluated under the identical entity-matching protocol: a \textbf{BERT-base clinical NER} model (\nolinkurl{samrawal/bert-base-uncased_clinical-ner}, 110M parameters) and \textbf{GatorTron-base clinical NER} (\nolinkurl{longluu/Clinical-NER-MedMentions-GatorTronBase}, a clinical-NER fine-tune of the 345M-parameter UF/NVIDIA GatorTron-base encoder on MedMentions). On mention-level extraction these reach F1 0.875 / 0.880 / 0.888 (BERT-base) and 0.840 / 0.847 / 0.876 (GatorTron) on 4CE / CORAL--Breast / CORAL--Pancreas, versus CLINES (o3-mini-medium) 0.906 / 0.881 / 0.913 (Figure~4D). Crucially, these encoder baselines emit entity spans only: they perform \emph{no} UMLS normalization, assertion, value+unit, or date extraction, so their code-level and attribute F1 are undefined (shown as 0 in Figure~4D). This makes explicit that the comparison is not span detection alone---where the encoders are competitive---but end-to-end normalized structuring, which they do not provide. We evaluate these models as released (no fine-tuning on our corpora), so the comparison reflects off-the-shelf clinical-NER capability rather than any tuning advantage on our part.}

\subsubsection*{R3.9 --- Figure 3 readability}
\reviewer{``Figure 3 would benefit from improvement in readability.''}
\response{Figure~3 has been redesigned with higher resolution (300 DPI), larger font sizes for axis labels and values, and clearer color separation between conditions. We also added per-bin sample sizes to Figure~3E (cf.\ R5.3.5).}

\subsubsection*{R3.10 --- Results section interpretation depth}
\reviewer{``The Results section would benefit from more descriptive and insightful interpretation.''}
\response{Each results subsection now opens with a one-sentence takeaway and closes with a brief mechanistic interpretation. The revised Results sections prioritize clinical meaning over numerical recitation.}

\subsubsection*{R3.11 --- Limitations underdeveloped}
\reviewer{``\dots\ the discussion remains underdeveloped. In addition to identifying limitations, the paper should more clearly outline possible paths for overcoming them.''}
\response{The Limitations subsection (the Discussion) has been expanded substantially. Each limitation is now paired with a specific path for future work: e.g., (i) prompt sensitivity $\rightarrow$ automatic prompt-search and optimization workflows; (iii) hallucination $\rightarrow$ retrieval-grounded extraction with citation back to source spans; (iv) IAA $\rightarrow$ extension to full double-annotation on a larger validation cohort.}

\subsubsection*{R3.12 --- Missing Conclusion section}
\reviewer{``Current writing lacks a Conclusion section which is supposed to contain the main take-home message.''}
\response{A new \textbf{Conclusion} section has been added at the end of the Discussion, summarizing the main contribution, the cross-cutting empirical findings (IAA bounds, hallucination breakdown, baseline-relative gains), and the recommended path to deployment.}

\subsubsection*{R3.13 --- Future Work too brief}
\reviewer{``The Future Work section would benefit from greater elaboration.''}
\response{Expanded. The Future Work subsection now covers: (i) domain-adaptive fine-tuning targets (oncology subspecialties, pediatrics, ICU); (ii) adaptive feedback workflows (online refinement of prompts from clinician corrections); (iii) multilingual extensions (Spanish, Mandarin); (iv) retrieval-augmented hallucination mitigation with span-level provenance; (v) integration with OMOP-CDM and FHIR (cf.\ R5.4.4).}

\hrulefill

\section{Reviewer 4}

\subsubsection*{R4.1 --- Chunking strategy detail}
\response{See R3.3.}

\subsubsection*{R4.2 --- Step 4 cross-chunk consistency}
\response{See R3.4.}

\subsubsection*{R4.3 --- Error analysis}
\response{See R1.9. Seven-class hallucination taxonomy plus case studies cover the error-categorization request.}

\subsubsection*{R4.4 --- Cost analysis for proprietary models}
\response{See R1.5. Full cost/latency/token-count table per dataset per model added (Figure~4B).}

\subsubsection*{R4.5 --- Hybrid LLM (different models per stage)}
\reviewer{``\dots\ it would be interesting to explore whether different LLMs could be used at different stages of the pipeline. For example, LLaMA could be used for earlier stages such as entity extraction, while more powerful models could be reserved for steps requiring deeper reasoning.''}
\response{We thank the reviewer for this interesting suggestion. The current modular design already supports per-stage model selection (the LLM provider is configurable per stage), but a systematic empirical evaluation of hybrid configurations is beyond the scope of the current revision. We have added this explicitly as a future-work item (cf.\ R3.13).}

\hrulefill

\section{Reviewer 5}

\subsection*{2. General considerations and formatting}

\subsubsection*{R5.2.1 --- Abstract structure}
\response{See R1.18. Abstract reframed to BMJ DH\&AI standard headings.}

\subsubsection*{R5.2.2 --- Summary box}
\response{Added immediately after the abstract: ``What is already known''; ``What this study adds''; ``How this study may affect research, practice, or policy''.}

\subsubsection*{R5.2.3 --- CLINES acronym inconsistency}
\response{Unified throughout: ``CLINES: \textbf{C}linical \textbf{L}LM-based \textbf{I}nformation \textbf{E}xtraction and \textbf{S}tructuring Agent''. The title and the page~6 expansion now match exactly.}

\subsubsection*{R5.2.4 --- Graphics resolution}
\response{All figures regenerated at 300 DPI, with larger font sizes and clearer color separation. Figure~3 in particular has been redesigned to address legibility concerns; shaded confidence ribbons have higher opacity (cf.\ R3.9).}

\subsection*{3. Minor findings, suggestions for improvement}

\subsubsection*{R5.3.1 --- Prompt suite availability and URL}
\response{We confirm that the full prompt suite, with all few-shot examples, will be released alongside an inference-only public reference implementation at the URL specified in the revised Data Availability statement. The implementation is licensed under the MIT License and provides single-note inference using Azure OpenAI--compatible endpoints (user-supplied credentials).}

\subsubsection*{R5.3.2 --- Quantitative hallucination analysis}
\response{See R1.9. Seven-class taxonomy and case studies added.}

\subsubsection*{R5.3.3 --- STROBE inappropriate}
\response{See R1.4. STROBE replaced by MI-CLAIM and TRIPOD+AI.}

\subsubsection*{R5.3.4 --- SemChunk and clinical-reasoning continuity}
\reviewer{``Clinical reasoning is not constrained by the semantic units kept by SemChunk\dots\ a missed reasoning across boundaries will still be lost to extraction\dots\ There is no detailed description of the context window mechanism\dots''}
\response{Acknowledged. The revised Methods clarifies that SemChunk is used as a \emph{boundary-detection heuristic} to align chunks with sentence/paragraph breaks, not as a guarantee of clinical-reasoning continuity. We additionally describe the cross-chunk reconciliation mechanism (Step~4; cf.\ R3.4) which joins chunk-local outputs by global entity tag and metadata anchors; because chunks are non-overlapping, the reconciler does not perform cross-chunk span deduplication, and same-surface mentions in different chunks remain separate records (preserving their chunk-local context's attribute values). We do not claim that semantic chunking eliminates all cross-chunk reasoning loss; this is now explicitly listed as a limitation (the Discussion).}

\subsubsection*{R5.3.5 --- Notes-per-bin in Figure 3E}
\response{Acknowledged. The total annotated document counts per corpus ($n=13$ to $n=24$) are stated in the revised Methods (``Study design and data''), and Supplementary Materials (``Extended Results: Statistical confidence and stability under small evaluation sets'') makes the small-sample uncertainty explicit: on CORAL--Breast ($n=13$) the bootstrap standard deviation of F1 reaches $\approx 0.04$--$0.05$, so F1 differences below $\approx 0.10$ are not statistically resolvable. The revised Discussion warns against over-interpreting the smallest length-bins, and the Figure~3E length-trend visualisation itself was retained without alteration so as not to obscure its primary message (stable CLINES performance across the length range).}

\subsection*{4. Major issues}

\subsubsection*{R5.4.1 --- Single-prompt LLM baseline insufficient}
\reviewer{``\dots\ The most critical comparison (o3-mini single-prompt versus o3-mini in CLINES) is entirely absent\dots\ The clinical NLP prompting literature from 2024–2025 has established that modern LLMs perform substantially better\dots\ with chain-of-thought instructions, few-shot clinical examples, and explicit task decomposition. None of these are present in the CLINES baseline.''}
\response{We have added both requested baseline conditions (cf.\ R1.7): (i) o3-mini single-prompt and (ii) GPT-4o chain-of-thought single-prompt with structured-schema instruction. The previously-implied claim of architectural innovation as the sole driver of the F1 gains is reframed in the revised Discussion to acknowledge that the observed gains reflect the \emph{combined} contributions of backbone model, prompt engineering, reasoning budget, and pipeline decomposition. The component-wise ablation (Figure~4C, cf.\ R1.8) further disentangles these contributions.}

\subsubsection*{R5.4.2 --- IAA Cohen's $\kappa$ requirement}
\response{See R1.1 response. Mention F1 (0.82), assertion Cohen's $\kappa$ (0.85), and per-attribute exact-match rates reported on 3 notes per dataset (10--29\% per-dataset coverage); per-note breakdown in Supplementary Table~S5. MIMIC-III F1 reframed as agreement-with-annotator. IAA reporting limitation now part of main Methods.}

\subsubsection*{R5.4.3 --- Scalability and deployment-condition claims}
\reviewer{``CLINES takes clean, de-identified, plain-text notes as input. The manuscript is entirely silent on the processes required to produce that input operationally\dots\ The only single 8-GPU H100 configuration that accommodates this model is FP8 quantisation\dots\ The GPT-4o and o3-mini variants require external API calls incompatible with most hospital governance frameworks and, in European jurisdictions, with GDPR.''}
\response{The revised main-text Discussion adds a brief ``Deployment, interoperability, and ethical considerations'' subsection; the full deployment discussion is in Supplementary Materials (``Deployment, Interoperability, and Ethical Considerations''). The added content:
\begin{itemize}[leftmargin=*]
    \item Specifies that the Llama-3.1-405B variant was run in \textbf{FP8 precision} on $8\times$NVIDIA H100 GPUs, and acknowledges the quality trade-off relative to FP16.
    \item Explicitly enumerates the preprocessing steps required to produce de-identified plain-text input from raw EHR sources (notably PHI de-identification, format normalization, and section/document segmentation), and acknowledges that this preprocessing is institutionally heterogeneous and not feasible for all settings.
    \item Discusses governance constraints for the commercial-API variants, including GDPR considerations in European jurisdictions, and notes that the Microsoft Azure OpenAI deployment used in this work (HIPAA BAA-covered) does not satisfy GDPR data-residency requirements outside the EU region. We recommend the open-weight backbone (Llama-3.1-405B or its successors) for institutions with EU data-residency requirements.
\end{itemize}}
\fromms{Supplementary Materials, ``Deployment, Interoperability, and Ethical Considerations'' (Deployment considerations); brief pointer in main-text Discussion (``Deployment, interoperability, and ethical considerations'')}{``\textcolor{blue}{the open-weight Llama-3.1-405B backbone was served in FP8 precision on 8$\times$H100 GPUs---the only single-node configuration that accommodates a 405B model---which entails a quality trade-off relative to FP16 and an infrastructure footprint beyond the reach of many small and medium centers\dots\ this configuration does not by itself satisfy GDPR data-residency requirements outside the relevant region, so institutions with EU data-residency obligations should prefer a locally deployed open-weight backbone.}''}

\subsubsection*{R5.4.4 --- UMLS CUI vs SNOMED-CT, LOINC, RxNorm, OMOP-CDM}
\reviewer{``CLINES produces UMLS CUIs concatenated with preferred term strings\dots\ UMLS's knowledge structure is explicitly not compliant with FHIR or any known healthcare interoperability standard; production FHIR resources require SNOMED-CT, LOINC, or RxNorm codes, not raw CUIs. OMOP-CDM\dots\ is unaddressed.''}
\response{Tempered. Supplementary Materials (``Deployment, Interoperability, and Ethical Considerations'', subsection ``Interoperability with production terminologies'') now enumerates the intermediate mapping steps required to move from UMLS CUI output to SNOMED-CT (via UMLS's cross-source mapping), LOINC (for laboratory observables), RxNorm (for medications), or OMOP-CDM (which natively supports UMLS-vocabulary ingestion via the OMOP Vocabularies, but requires standardization of attribute fields). We have downgraded the previous ``FHIR-ready'' implication and replaced it with ``ontology-grounded intermediate representation suitable for downstream mapping into production terminologies including SNOMED-CT, LOINC, RxNorm, and OMOP-CDM.'' The main-text Discussion subsection ``Deployment, interoperability, and ethical considerations'' includes a brief pointer; Supplementary ``Mapping UMLS Concepts to Source Terminologies'' gives a worked crosswalk table.}
\fromms{Supplementary Materials, ``Deployment, Interoperability, and Ethical Considerations'' (Interoperability with production terminologies); brief pointer in main-text Discussion}{``\textcolor{blue}{CLINES emits UMLS CUIs paired with preferred terms. UMLS is an ontology-grounded intermediate representation rather than a production interchange format: downstream use in FHIR resources or common data models requires mapping CUIs onward to the appropriate target vocabularies---SNOMED~CT (via UMLS cross-source mappings), LOINC\dots\ or RxNorm\dots---and OMOP-CDM ingestion via its standardized vocabularies\dots\ we recommend retaining a SNOMED~CT intermediate layer in deployments that target FHIR.}''}

\subsubsection*{R5.4.5 --- Omission of contemporary state-of-the-art comparators (Healthcare NLP, CLEAR, BMJ paper)}
\reviewer{``a) John Snow Labs Healthcare NLP / Spark NLP\dots\ b) CLEAR\dots\ npj Digit. Med. 8, 45 (2025)\dots\ c) Ntinopoulos V, et al.\ BMJ Health \& Care Informatics 2025.''}
\response{We thank the reviewer for these pointed references. Each is now discussed in the main-text Discussion (``Positioning versus existing approaches'') with the bulk of the matched-comparison reasoning in Supplementary Materials (``Extended Results: Positioning versus contemporary production and research systems''), and the Ntinopoulos output-consistency check is reported as its own Supplementary section.
\begin{itemize}[leftmargin=*]
    \item \textbf{John Snow Labs Healthcare NLP / Spark NLP}: a large, separately-maintained production clinical-NLP platform that implements many capabilities CLINES describes. We now position CLINES as a \emph{transparent, fully-specified open reference implementation} prioritizing reproducibility of the prompt suite, complementary to (rather than competing with) production suites. We did not include a head-to-head re-benchmark: a fair matched comparison (identical preprocessing, identical gold standard, identical span-matching protocol) against a continuously-updated commercial platform is a substantial undertaking beyond this revision's scope, and we now state this absence explicitly as a limitation rather than implying equivalence.
    \item \textbf{CLEAR (Lopez et al., npj Digit.\ Med.\ 2025)}: discussed in the Supplementary positioning subsection as a directly comparable LLM-augmented pipeline. While we did not re-implement CLEAR for direct empirical comparison, we contrast the architectural choices (CLEAR's retrieval-augmented generation focus vs CLINES's normalization-first multi-step decomposition).
    \item \textbf{Ntinopoulos et al., BMJ HCI 2025}: cited in the Supplementary positioning subsection. We note that the 18-model evaluation included output consistency assessment which the original CLINES submission did not, and we accordingly added an output-consistency robustness check to Supplementary Materials (``Output Consistency Across Repeated Runs''). Across $R=5$ independent runs of the full pipeline on a fixed 20-note subset (all runs completed; no missing outputs), matching mentions by semantic equivalence (shared UMLS code or shared canonical surface form, absorbing span-boundary jitter and surface-form synonyms), pairwise mention-set Jaccard averaged $0.752$ and code-set Jaccard $0.622$. For the $27{,}605$ (mention-class, field) cells with $\geq 2$ contributing runs, the modal field value was matched in $89.9\%$ of contributing runs on average (unit $94.9\%$, end date $94.2\%$, begin date $91.3\%$, assertion status $84.9\%$, value $84.3\%$). The pattern---roughly three quarters of mentions are semantically shared between any two runs, and per-field content of recovered mentions is highly stable---supports majority-vote aggregation across $R\geq 3$ runs when reproducibility matters and is now documented as a deployment recommendation.
\end{itemize}}

\hrulefill

\section*{Closing}

We thank the Editor and the reviewers again for the depth and constructiveness of their feedback, which has materially strengthened the manuscript. We hope the revised version addresses their concerns and look forward to further consideration.

\end{document}
